% !TeX root = prueba02_version02.tex
% !TeX program = pdflatex
% Prueba 2: versión 2, dificultad media, 120 minutos, 10 puntos.
\documentclass[10pt,a5paper]{article}
\newif\ifmostrarSoluciones
\mostrarSolucionesfalse
\usepackage{estilo_apuntes}
\dosPorCarillafalse % Use \dosPorCarillatrue para imprimir dos copias por carilla.
\encabezadoprueba{Prueba 2}{Jueves 12 de noviembre de 2026}

\begin{document}
\begin{center}
  {\color{Azul}\large\bfseries Prueba 2}
\end{center}

\noindent\textbf{Nombre:} \hrulefill

\medskip
\noindent\textbf{Indicaciones}
\begin{itemize}
  \item La duración de la prueba es de 1 hora y 30 minutos. Se califica sobre 10 puntos.
  \item Escriba el desarrollo de manera clara y ordenada. Justifique las respuestas e indique las propiedades y restricciones utilizadas. Al aplicar un teorema, verifique sus hipótesis y exprese su conclusión. Una respuesta sin justificación recibirá únicamente crédito parcial.
  \item Resuelva con las herramientas de límites y continuidad estudiadas, sin derivadas ni regla de L'Hôpital. No se solicitan demostraciones de teoremas.
  \item No se permite el uso de dispositivos electrónicos. Cualquier intento de copia penalizará la evaluación con una calificación total de cero.
\end{itemize}

\noindent\textbf{Ejercicios}
\begin{enumerate}
  \item \textbf{(4 puntos)} Considere
  \[
    R(x)=\frac{x+1}{(x-2)^2},\qquad
    D_R=\mathbb{R}\setminus\{2\},
  \]
  y
  \[
    H(x)=\sqrt{x^2+4x+8}-x,\qquad D_H=\mathbb{R}.
  \]
  \begin{enumerate}
    \item \textbf{(2 puntos)} Calcule los límites laterales de $R$ en $2$ y determine si se puede expresar su comportamiento mediante un único límite bilateral infinito. Determine todas sus asíntotas verticales.
    \item \textbf{(2 puntos)} Calcule los límites de $H$ cuando $x\to+\infty$ y $x\to-\infty$. Determine sus asíntotas horizontales e indique a qué extremo corresponde cada una.
  \end{enumerate}

  \item \textbf{(3 puntos)} Para $a,b\in\mathbb{R}$, sea
  \[
    F_{a,b}\colon[-3,\infty)\to\mathbb{R},\qquad
    F_{a,b}(x)=\begin{cases}
      \dfrac{\sqrt{x+3}-2}{x-1},&-3\leq x<1,\\[2mm]
      a,&x=1,\\
      \dfrac{bx+1}{x+1},&x>1.
    \end{cases}
  \]
  \begin{enumerate}
    \item \textbf{(2 puntos)} Determine $a,b$ para que $F_{a,b}$ sea continua en todo su dominio. Justifique el comportamiento en el punto de unión, en el extremo $-3$ y en el interior de cada tramo.
    \item \textbf{(1 punto)} Ahora tome $a=0$, $b=-1/2$. Clasifique la discontinuidad en $x=1$ e indique si se puede eliminar cambiando únicamente $F_{0,-1/2}(1)$. Si es posible, especifique el valor que debe asignarse.
  \end{enumerate}
  % Aplicación concreta de máximo--mínimo y valor intermedio.
  \item \textbf{(3 puntos)} Considere
  \[
    g(x)=x+\frac4x,\qquad x\in[1,4].
  \]
  \begin{enumerate}
    \item \textbf{(1 punto)} Justifique que $g$ alcanza un máximo y un mínimo absolutos en $[1,4]$, indicando el teorema y verificando sus hipótesis.
    \item \textbf{(1 punto)} Determine el valor mínimo y el punto donde se alcanza, sin utilizar derivadas. Puede estudiar el signo de $g(x)-4$.
    \item \textbf{(1 punto)} Justifique mediante el teorema del valor intermedio que la ecuación
    \[
      x+\frac4x=\frac92
    \]
    tiene al menos dos soluciones distintas en $(1,4)$, una a cada lado de $x=2$. No se pide calcularlas.
  \end{enumerate}
\end{enumerate}

\ifmostrarSoluciones
\newpage
\begin{solucion}[de la versión 2]
\begin{enumerate}
  \item
  \begin{enumerate}
    \item Cerca de $2$, el numerador $x+1$ es positivo y tiende a $3$. El denominador $(x-2)^2$ tiende a $0$ mediante valores positivos desde ambos lados. Por tanto,
    \[
      \lim_{x\to2^-}R(x)=+\infty,\qquad
      \lim_{x\to2^+}R(x)=+\infty.
    \]
    Se puede escribir $\lim_{x\to2}R(x)=+\infty$. La única asíntota vertical es $x=2$, pues en los demás puntos el denominador no se anula y el límite es finito.
    \item El radicando es $(x+2)^2+4>0$ para todo $x$. Para estudiar $x\to+\infty$, racionalizamos para $x>0$:
    \begin{align*}
      H(x)&=\frac{4x+8}{\sqrt{x^2+4x+8}+x}\\
          &=\frac{4+8/x}{\sqrt{1+4/x+8/x^2}+1}.
    \end{align*}
    En el segundo paso se usa $\sqrt{x^2}=x$, válido para $x>0$. Así,
    \[
      \lim_{x\to+\infty}H(x)=\frac4{1+1}=2.
    \]
    Si $x\to-\infty$, la raíz es no negativa y $H(x)\geq-x\to+\infty$. Por comparación,
    \[
      \lim_{x\to-\infty}H(x)=+\infty.
    \]
    La única asíntota horizontal es $y=2$, hacia $+\infty$. No hay asíntota horizontal hacia $-\infty$.
  \end{enumerate}

  \item
  \begin{enumerate}
    \item Para $-3\leq x<1$, al multiplicar por el conjugado se obtiene
    \[
      \frac{\sqrt{x+3}-2}{x-1}=\frac1{\sqrt{x+3}+2}.
    \]
    Por ello, el límite izquierdo en $1$ es $1/4$. El límite derecho es $(b+1)/2$, y el valor de la función es $a$. La continuidad exige
    \[
      \frac14=a=\frac{b+1}{2},
    \]
    de donde $a=1/4$ y $b=-1/2$. En los interiores, las expresiones son continuas y sus denominadores no se anulan. En el extremo $-3$, la expresión simplificada da
    \[
      \lim_{x\to-3^+}F_{a,b}(x)=\frac12=F_{a,b}(-3),
    \]
    por lo que también hay continuidad lateral en el extremo. Así, con los valores encontrados, la función es continua en todo $[-3,\infty)$.
    \item Si $a=0$ y $b=-1/2$, ambos límites laterales en $1$ valen $1/4$, pero $F_{0,-1/2}(1)=0$. La discontinuidad es removible: se elimina asignando el valor $1/4$ en $x=1$, sin modificar los demás valores de la función.
  \end{enumerate}
  \item
  \begin{enumerate}
    \item La función $g(x)=x+4/x$ es continua para $x>0$, en particular en todo $[1,4]$. Como el intervalo es cerrado y acotado, el teorema de máximo--mínimo garantiza que alcanza ambos extremos absolutos.
    \item Para $x\in[1,4]$,
    \[
      g(x)-4=x+\frac4x-4=\frac{(x-2)^2}{x}\geq0,
    \]
    porque el numerador es un cuadrado y el denominador es positivo. Por tanto, $g(x)\geq4$. La igualdad ocurre solamente en $x=2$, que pertenece al intervalo. El mínimo absoluto es $4$ y se alcanza únicamente en $x=2$.
    \item Se tiene $g(1)=5$, $g(2)=4$ y $g(4)=5$. La función es continua en $[1,2]$ y $[2,4]$, y $4<9/2<5$. El teorema del valor intermedio aplicado en cada intervalo garantiza un $c_1\in(1,2)$ y un $c_2\in(2,4)$ tales que
    \[
      g(c_1)=g(c_2)=\frac92.
    \]
    Las soluciones son distintas porque pertenecen a intervalos disjuntos. No es necesario resolver la ecuación para justificar su existencia.
  \end{enumerate}

\end{enumerate}
\end{solucion}
\fi

\end{document}
